\section{Methodology}
\label{sec:methodology}

\subsection{County-sticky edge weighting}

\subsubsection{Formulation}

Let $G = (V, E, w_V, w_E)$ be the census-tract adjacency graph from B.2,
where $w_E(u, v)$ is the shared boundary length between tracts $u$ and $v$.
Let $c : V \to \mathcal{C}$ map each tract to its county FIPS code.

The county-sticky edge weight is:
\[
  w_E'(u, v) = w_E(u, v) \cdot
  \begin{cases}
    \alpha_c & \text{if } c(u) = c(v) \quad \text{(intra-county edge)} \\
    1        & \text{if } c(u) \neq c(v) \quad \text{(inter-county edge)}
  \end{cases}
\]
where $\alpha_c \geq 1$ is the county stickiness parameter.

\paragraph{Interpretation.}
When $\alpha_c = 1$, the weighting reduces to the B.2 boundary-length
baseline.  When $\alpha_c > 1$, intra-county cuts are penalised relative
to inter-county cuts: the algorithm prefers to place district boundaries
along county lines rather than through county interiors.
When $\alpha_c \to \infty$, the formulation approaches the hard-constraint
limit (no intra-county cuts), which is infeasible for large counties.

\subsubsection{Planarity and connectivity}

The county-sticky weighting preserves all structural properties of the
adjacency graph: planarity (edge set is unchanged), connectivity (no edges
are removed), and the positive edge-weight requirement for METIS.
Correctness of METIS under arbitrary positive edge weights follows from
Karypis and Kumar~\citep{karypis1998}.

\subsubsection{County split metrics}

We define:
\begin{itemize}
  \item \emph{County split incidences}: for each county $c$, let
    $m_c$ be the number of districts intersecting $c$.
    The split count is $\sum_c \max(0, m_c - 1)$.
    This counts the number of extra district pieces created beyond the
    first district touching each county.
  \item \emph{Multi-county districts}: the number of districts that
    draw tracts from more than one county.
  \item \emph{Perfect county preservation}: a district that contains
    exactly the tracts of one or more complete counties (no partial
    counties).
\end{itemize}

The first two metrics are related but not interchangeable.
County split incidences are county-centered: they ask how many additional
district pieces counties are broken into.
Multi-county districts are district-centered: they ask how many districts cross
at least one county line.
A single district can contribute to multiple county split incidences, and a
single split county can affect multiple districts.
Thus the national headline numbers answer different questions: 487 $\to$ 323
reports fewer extra county pieces, while 312 $\to$ 198 reports fewer districts
that span county boundaries.

\subsection{Parameter sweep design}

We sweep $\alpha_c \in \{1.0, 1.5, 2.0, 2.5, 3.0, 3.5, 5.0, 10.0\}$ on all 50
states using the 2020 Census congressional maps ($k$ given by 2020
apportionment).  For each $(\text{state}, \alpha_c)$ pair:
\begin{enumerate}
  \item Build the county-sticky adjacency graph using the tract FIPS-to-county
    mapping from TIGER/Line data~\citep{tiger-census}.
  \item Run METIS recursive bisection with fixed seed ($s_0$, DIA formula).
  \item Record: mean Polsby--Popper ($\overline{PP}$), county splits,
    multi-county districts, and max population deviation.
\end{enumerate}

Total runs: $50 \text{ states} \times 8 \text{ } \alpha_c \text{ values} =
400$ redistricting runs.

\subsection{Pareto frontier analysis}

We identify the Pareto frontier of (county splits, compactness loss) by
plotting county split reduction vs.\ $\overline{PP}$ loss for each
$\alpha_c$ value.  The ``elbow'' of the Pareto frontier identifies the
$\alpha_c$ value beyond which further increases produce diminishing
county-split reduction at disproportionate compactness cost.

In the atlas framing, the candidates are not alternative district maps alone
but alternative weighted graph policies.
Each $\alpha_c$ candidate produces a graph, a deterministic plan, and a
measurement record containing county splits, multi-county districts,
compactness, and population deviation.
The decision rule is the Pareto elbow: select the smallest weight that removes
avoidable county splits without paying a disproportionate compactness cost.
The evidence path must preserve the whole sweep, because a single selected map
cannot show which county splits were unavoidable, which were avoided by the
weighting policy, or where the trade-off became too expensive.

\subsection{State selection for case studies}

We provide detailed case studies for:
\begin{itemize}
  \item \textbf{Iowa} ($k{=}4$, 99 counties): many small counties relative
    to district count; county preservation is both achievable and
    constitutionally required.
  \item \textbf{Texas} ($k{=}38$, 254 counties): many counties of
    heterogeneous size; some counties require splitting; tests the algorithm's
    behaviour at the large-$k$ limit.
  \item \textbf{California} ($k{=}52$, 58 counties): large counties that
    \emph{must} be split; tests whether the algorithm avoids unnecessary splits.
\end{itemize}
